\section{实轴上的小波}
\label{sec:wavelet-real-line}

区间上的Haar分解由两部分配合：区间指示函数记录平均水平，左右两半取相反值的函数记录细节。这两类函数都由一个固定的函数平移、缩放而来。在全实轴上，位置$k$不再限于$0\leq k<2^j$，而要取遍全部整数，才能覆盖整个定义域。本节先把Haar的结构抽象成正交多分辨率构造，再在这一结构下写出$L^2(\mathbb R)$中的展开，最后讨论为什么要用比Haar更光滑、能抵消更多项的小波。

本节只考虑实值的尺度函数和小波，内积$\langle f,g\rangle=\int_{\mathbb R}\overline{f(x)}g(x)\,dx$。对函数$u$和整数$j,k$，记
\begin{equation}
u_{j,k}(x)=2^{j/2}u(2^jx-k).
\label{eq:wavelet-dilate-translate}
\end{equation}
变量代换$y=2^jx-k$给出$\int_{\mathbb R}\lvert u_{j,k}\rvert^2\,dx=\int_{\mathbb R}\lvert u\rvert^2\,dy$，所以因子$2^{j/2}$使缩放不改变$L^2$范数；同理$\langle u_{j,k},v_{j,k'}\rangle=\langle u_{0,k},v_{0,k'}\rangle$，同一层内的内积不依赖$j$。

\subsection{尺度函数与母小波}
\label{sec:wavelet-scaling-mother}

先把Haar的情形写到实轴上。令$\varphi=\mathbf{1}_{[0,1)}$，$\psi=\mathbf{1}_{[0,1/2)}-\mathbf{1}_{[1/2,1)}$，并令$V_j$为$L^2(\mathbb R)$中在每个区间$[k2^{-j},(k+1)2^{-j})$（$k\in\mathbb Z$）上取常值的函数组成的闭子空间，这里$j$可以是任意整数。与区间上一样，$\{\varphi_{j,k}\}_{k\in\mathbb Z}$是$V_j$的标准正交基，$V_j\subset V_{j+1}$。不同之处在于每层有无限多个位置，并且$j$也可以取负数，对应长度为$2^{-j}>1$的粗区间。

\paragraph{多分辨率的条件}
Haar的例子满足以下几条，它们正是一般构造需要的结构。

\begin{definition}[正交多分辨率分析]
\label{def:wavelet-mra}
$L^2(\mathbb R)$中的一列闭子空间$\{V_j\}_{j\in\mathbb Z}$连同函数$\varphi\in V_0$，若满足
\begin{enumerate}
\item 嵌套：$V_j\subset V_{j+1}$对一切$j\in\mathbb Z$成立；
\item 缩放：$f\in V_j$当且仅当$f(2\,\cdot)\in V_{j+1}$；
\item 平移与基：$\{\varphi(\cdot-k)\}_{k\in\mathbb Z}$是$V_0$的标准正交基；
\item 稠密与分离：$\bigcup_jV_j$在$L^2(\mathbb R)$中稠密，且$\bigcap_jV_j=\{0\}$，
\end{enumerate}
则称它为一个\term{正交多分辨率分析}（orthogonal multiresolution analysis），$\varphi$称为\term{尺度函数}（scaling function）。
\end{definition}

由缩放与平移两条，$\{\varphi_{j,k}\}_{k\in\mathbb Z}$是$V_j$的标准正交基：$f\in V_j$当且仅当$f(2^{-j}\cdot)\in V_0$，而式~\eqref{eq:wavelet-dilate-translate}的归一化保持标准正交。对Haar，稠密性可以仿照定理~\ref{thm:wavelet-dyadic-dense}证明：先用截断把$f$换成支撑在某个有界区间上的函数，再在这个区间上重复四步逼近。分离条件也可以直接验证：若$f\in V_{-m}$对一切$m\geq0$成立，则$f$在$[0,2^m)$与$[-2^m,0)$上都取常值，从而在两条半轴上各为常数，平方可积迫使这两个常数为零。

\paragraph{二尺度关系}
相邻两层由一个系数序列连接。$\varphi\in V_0\subset V_1$，而$\{\varphi_{1,k}\}_{k\in\mathbb Z}$是$V_1$的标准正交基，所以
\begin{equation}
\varphi(x)=\sum_{k\in\mathbb Z}h_k\varphi_{1,k}(x)=\sqrt2\sum_{k\in\mathbb Z}h_k\varphi(2x-k),\qquad h_k=\langle\varphi_{1,k},\varphi\rangle,
\label{eq:wavelet-two-scale}
\end{equation}
级数在$L^2$中收敛，且$\sum_k\lvert h_k\rvert^2=\lVert\varphi\rVert_2^2=1$。式~\eqref{eq:wavelet-two-scale}称为\term{二尺度关系}（two-scale relation）：尺度函数等于自身压缩一半后的若干平移的组合。Haar的情形是$h_0=h_1=1/\sqrt2$，其余为零，式~\eqref{eq:wavelet-two-scale}就是$\mathbf{1}_{[0,1)}=\mathbf{1}_{[0,1/2)}+\mathbf{1}_{[1/2,1)}$。

\paragraph{母小波}
细节空间仍定义为正交补：$W_j=V_{j+1}\cap V_j^{\perp}$，于是$V_{j+1}=V_j\oplus W_j$。我们希望$W_0$也由单个函数的整数平移张成。由二尺度系数可以写出这个函数。

\begin{theorem}[由多分辨率构造小波]
\label{thm:wavelet-mother-from-mra}
设$\{V_j\},\varphi$是实值的正交多分辨率分析，$h_k$由式~\eqref{eq:wavelet-two-scale}给出。令
\begin{equation}
\psi(x)=\sqrt2\sum_{k\in\mathbb Z}g_k\varphi(2x-k),\qquad g_k=(-1)^kh_{1-k}.
\label{eq:wavelet-mother}
\end{equation}
则$\{\psi(\cdot-k)\}_{k\in\mathbb Z}$是$W_0$的标准正交基；从而对每个整数$j$，$\{\psi_{j,k}\}_{k\in\mathbb Z}$是$W_j$的标准正交基。
\end{theorem}

称这样的$\psi$为\term{母小波}（mother wavelet），称$\psi_{j,k}$为\term{小波}（wavelet）。这里不给出定理的完整证明。标准证明借助第~\ref{chap:fourier-transform}章的Fourier变换，把“整数平移标准正交”改写成频域上的周期恒等式，再在频域中验证$\psi$的平移既与$V_0$正交又张成整个$W_0$。

可以直接检验其中最简单的一步：$\psi$与$\varphi$本身正交。两者都在$V_1$中，用标准正交基$\{\varphi_{1,k}\}$的系数计算，
\[
\langle\varphi,\psi\rangle=\sum_kh_kg_k=\sum_k(-1)^kh_kh_{1-k}.
\]
把下标$k$与$1-k$配对，两项分别是$(-1)^kh_kh_{1-k}$和$(-1)^{1-k}h_{1-k}h_k$，由$(-1)^{1-k}=-(-1)^k$恰好抵消，所以和为零。对Haar，$g_0=h_1=1/\sqrt2$，$g_1=-h_0=-1/\sqrt2$，式~\eqref{eq:wavelet-mother}给出$\psi(x)=\varphi(2x)-\varphi(2x-1)$，正是Haar函数。

\paragraph{不是任意的平移缩放都构成小波基}
定义~\ref{def:wavelet-mra}的每一条都有实际作用。取$u=\mathbf{1}_{[0,2)}/\sqrt2$，则$\lVert u\rVert_2=1$，但$\langle u,u(\cdot-1)\rangle=\int_1^2\frac12\,dx=\frac12\neq0$，它的整数平移不正交，不满足第3条。再取$u=\mathbf{1}_{[0,1)}$充当“母小波”：同一层内的平移正交，跨层却不正交，$\langle u_{0,0},u_{1,0}\rangle=\int_0^{1/2}\sqrt2\,dx=\sqrt2/2$。所以只把某个函数作平移和缩放，并不自动得到标准正交基；小波的正交性来自多分辨率结构和式~\eqref{eq:wavelet-mother}的特殊选择。

\subsection{实轴上的小波展开}
\label{sec:wavelet-real-line-expansion}

在已经有了正交多分辨率分析和定理~\ref{thm:wavelet-mother-from-mra}的母小波之后，展开的建立与区间上完全平行。

\begin{theorem}[实轴上的小波展开]
\label{thm:wavelet-real-line-expansion}
设$\{V_j\},\varphi$是正交多分辨率分析，$\psi$由定理~\ref{thm:wavelet-mother-from-mra}给出。固定整数$j_0$，则
\[
\{\varphi_{j_0,k}:k\in\mathbb Z\}\cup\{\psi_{j,k}:j\geq j_0,\ k\in\mathbb Z\}
\]
是$L^2(\mathbb R)$的标准正交基。对每个$f\in L^2(\mathbb R)$，在$L^2$范数下
\[
f=\sum_{k\in\mathbb Z}\langle\varphi_{j_0,k},f\rangle\varphi_{j_0,k}+\sum_{j\geq j_0}\sum_{k\in\mathbb Z}\langle\psi_{j,k},f\rangle\psi_{j,k},
\]
且系数的平方和等于$\lVert f\rVert_2^2$。
\end{theorem}

\begin{proof}
由$V_{j+1}=V_j\oplus W_j$反复分解，对$J>j_0$有$V_J=V_{j_0}\oplus W_{j_0}\oplus\cdots\oplus W_{J-1}$，各项两两正交。由定理~\ref{thm:wavelet-mother-from-mra}，每个$W_j$有标准正交基$\{\psi_{j,k}\}_k$，所以所给函数族是标准正交的，其闭线性张成包含每个$V_J$。由稠密条件，它是完备的标准正交系；展开式与Parseval等式由第~\ref{chap:inner-product-spaces}章的正交展开得到。
\end{proof}

\paragraph{与区间展开的比较}
区间上每一层只有$2^j$个位置，实轴上每一层都有无限多个位置$k\in\mathbb Z$，所以同层内的求和本身也是一个级数，在$L^2$中收敛。若$f$与$\varphi,\psi$都有紧支撑，那么每一层只有有限个$\psi_{j,k}$的支撑与$f$的支撑相交，每层的非零系数仍是有限个，个数大致与$2^j$乘以$f$的支撑长度成正比。

\paragraph{实轴上粗项可以消失}
区间上的粗项不能省略，实轴上却可以让起始层$j_0\to-\infty$。记$P_j$为到$V_j$的正交投影。对$i<j$，$V_i\subset V_j$给出$P_if=P_iP_jf$，而$P_jf-P_if\perp P_if$，于是
\[
\lVert P_jf-P_if\rVert_2^2=\lVert P_jf\rVert_2^2-\lVert P_if\rVert_2^2 .
\]
$\lVert P_jf\rVert_2$随$j$减小而单调不增，有极限，上式说明$(P_jf)$在$j\to-\infty$时是Cauchy列，极限记为$g$。对每个固定的$i$，$j\leq i$时$P_jf\in V_i$，$V_i$闭，所以$g\in V_i$；于是$g\in\bigcap_iV_i=\{0\}$。这说明$\lVert P_{j_0}f\rVert_2\to0$（$j_0\to-\infty$），于是
\[
f=\sum_{j\in\mathbb Z}\sum_{k\in\mathbb Z}\langle\psi_{j,k},f\rangle\psi_{j,k},
\]
$\{\psi_{j,k}\}_{j,k\in\mathbb Z}$本身就是$L^2(\mathbb R)$的标准正交基。以Haar和$f=\mathbf{1}_{[0,1)}$为例，$m\geq0$时$P_{-m}f$在$[0,2^m)$上取平均值$2^{-m}$、在其余区间上为零，所以$\lVert P_{-m}f\rVert_2^2=2^{-2m}\cdot2^m=2^{-m}\to0$：同样的质量被摊到越来越长的区间上，平方积分趋于零。区间$[0,1]$上没有更长的区间可摊，常数函数的粗项始终保留。

\subsection{光滑性与消失矩}
\label{sec:wavelet-smoothness-moments}

Haar函数有两个明显的局限。其一，它有跳跃，所以任何有限项部分和都是阶梯函数，即使$f$本身很光滑，重构在每个分点上都有跳跃。其二，它只能抵消常数：$\int\psi=0$，但$\int_0^1x\psi(x)\,dx=1/8-3/8=-1/4\neq0$。这两点分别对应小波的两种性质。

\begin{definition}[光滑性与消失矩]
\label{def:wavelet-vanishing-moments}
设$\psi$是母小波。若$\psi$有$s$阶连续导数，称它具有$C^s$光滑性。设$r$是正整数，并且$\int_{\mathbb R}\lvert x\rvert^{r-1}\lvert\psi(x)\rvert\,dx<\infty$。若
\[
\int_{\mathbb R}x^m\psi(x)\,dx=0\qquad(m=0,1,\ldots,r-1),
\]
称$\psi$具有$r$个\term{消失矩}（vanishing moments）。
\end{definition}

可积性条件保证这些积分有意义，紧支撑的有界函数自动满足。由上面的计算，Haar小波恰有$1$个消失矩。存在紧支撑的正交小波，具有任意给定个数$r$的消失矩，其中最常用的一族是Daubechies小波，支撑长度为$2r-1$，随$r$增大光滑阶数也逐步提高；它们没有简单的闭式表达，由二尺度系数$h_k$确定，这里不给出构造。

\paragraph{消失矩让局部多项式不留痕迹}
有$r$个消失矩的小波，与任何次数小于$r$的多项式的内积为零：多项式$p(x)$经代换$y=2^jx-k$仍是$y$的次数小于$r$的多项式，逐项用消失矩即得$\langle\psi_{j,k},p\rangle=0$。所以若$f$在$\psi_{j,k}$的支撑附近很接近某个低次多项式，对应的系数就只反映$f$偏离这个多项式的那部分。下面的定理把这一点写成定量估计。

\begin{theorem}[消失矩与系数衰减]
\label{thm:wavelet-moment-decay}
设$\psi$实值，支撑在$[0,L]$中，具有$r$个消失矩。设$j,k$为整数，$f$为实值函数，在区间$S_{j,k}=[k2^{-j},(k+L)2^{-j}]$上有$r$阶导数且$\lvert f^{(r)}\rvert\leq M$。则
\[
\lvert\langle\psi_{j,k},f\rangle\rvert\leq\frac{ML^r}{r!}\,\lVert\psi\rVert_1\,2^{-j(r+1/2)} .
\]
\end{theorem}

\begin{proof}
$\psi_{j,k}$支撑在$S_{j,k}$上。在左端点$x_0=k2^{-j}$处对$f$作$r-1$阶Taylor展开，$f=p+R$，其中$p$是次数小于$r$的多项式；由Lagrange余项，$x\in S_{j,k}$时$\lvert R(x)\rvert\leq M\lvert x-x_0\rvert^r/r!\leq M(L2^{-j})^r/r!$。由消失矩，$\langle\psi_{j,k},p\rangle=0$，所以
\[
\lvert\langle\psi_{j,k},f\rangle\rvert=\lvert\langle\psi_{j,k},R\rangle\rvert\leq\sup_{S_{j,k}}\lvert R\rvert\cdot\lVert\psi_{j,k}\rVert_1 .
\]
最后由代换$y=2^jx-k$，$\lVert\psi_{j,k}\rVert_1=2^{j/2}\cdot2^{-j}\lVert\psi\rVert_1=2^{-j/2}\lVert\psi\rVert_1$。
\end{proof}

用Haar检验这个估计。$r=1$，$L=1$，$\lVert\psi\rVert_1=1$。对$f(x)=x$，$M=1$，定理给出$\lvert\langle\psi_{j,k},f\rangle\rvert\leq2^{-3j/2}$。直接计算：记$h=2^{-j}$，$a=kh$，则$\int_a^{a+h/2}x\,dx-\int_{a+h/2}^{a+h}x\,dx=-h^2/4$，所以$\langle\psi_{j,k},f\rangle=-2^{j/2}h^2/4=-2^{-3j/2}/4$，与定理给出的阶一致。若换成$f(x)=x^2$，同样的计算给出$\langle\psi_{j,k},f\rangle=-2^{j/2}(h^2/2)(a+h/2)$，仍是$2^{-3j/2}$的阶：Haar只有一个消失矩，抵消不了一次项，$f$再光滑也不能让系数衰减得更快。若换用紧支撑且具有两个消失矩的小波，一次项也被抵消，定理~\ref{thm:wavelet-moment-decay}取$r=2$、$M=2$，给出$x^2$的系数至多是$2^{-5j/2}$的阶。

\paragraph{两种性质的分工}
消失矩作用在被展开函数上：$f$在某处越光滑，那里的细节系数衰减越快；$f$在某处有跳跃，覆盖跳跃的少数系数就衰减得慢。光滑性作用在重构上：用光滑的小波作有限项部分和，得到的是光滑函数，用它逼近光滑的$f$时不会人为引入跳跃，在包含导数的范数下比较误差也才有意义。两者是对小波本身的要求，应与$f$的正则性区分开：前者是选择工具，后者是被研究对象。用小波系数的衰减刻画$f$的正则性时，所需的消失矩个数和光滑阶数要比被刻画的阶数更高，具体条件在第~\ref{chap:besov-spaces}章另行固定。
